\section{High Usability}
\label{sec:usability}

The design of \ours prioritizes usability, aiming to streamline the development process for multi-agent with LLMs and to ensure a smooth interaction experience for both users and developers. 
This section delves into how \ours flattens the learning curve and enhances the programmer's experience by introducing intuitive concepts and features that facilitate the creation of complex multi-agent applications.

\subsection{Syntactic Sugar for Multi-Agent Workflows}
Leveraging basic concepts introduced in Section \ref{sec:concept}, developers are empowered to construct sophisticated multi-agent applications. Nonetheless, directly coding each agent's message exchange can become cumbersome, as shown in \lstlistingname~\ref{lst:without_pipeline}. Recognizing this, \ours introduces two syntactic utilities: pipelines and message hubs, to abstract away the complexity and minimize repetition.

\begin{lstlisting}[language=Python, float=h, caption={Example of programming a sequential workflow with basic concepts in \ours.}, label={lst:without_pipeline}]
# set up agents: agent1 to agent5
# ...

msg = agent1(Msg ( " Alice " , " Hello ! " ))
msg = agent2(msg)
msg = agent3(msg)
msg = agent4(msg)
msg = agent5(msg)
\end{lstlisting}

\paragraph{Pipeline Abstraction} 
The pipeline abstraction reduces repetitive coding by encapsulating patterns of message transmission, including {\it sequential}, {\it conditional}, and {\it iterative} exchanges, into simple and reusable components. With these pipelines, developers can focus on the logic of agent interactions rather than the boilerplate code. \lstlistingname~\ref{lst:pipeline} illustrates how pipelines can be employed in both functional and object-oriented styles to create a clear and concise agent workflow.
Besides the sequential pipeline in the example, \ours also provides if-else, switch, while-loop, and for-loop pipelines, facilitating the programming of the multi-agent interactions.

\begin{lstlisting}[language=Python, float=h, caption={Using functional and object sequential pipeline to construct workflow in \ours.}, label={lst:pipeline}]
# set up agents: agent1 to agent5
# ...
from agentscope.pipelines import SequentialPipeline
from agentscope.pipelines.functional import sequentialpipeline

# using functional pipeline
x = sequentialpipeline([agent1, agent2, agent3, agent4, agent5], x)

# using object pipeline
pipe = SequentialPipeline([agent1, agent2, agent3, agent4, agent5])
x = pipe(x)
\end{lstlisting}  

\paragraph{Message Hub for Agent Communication}
In multi-agent systems, especially when integrated with LLMs, efficiently managing communication among a group of agents is essential. The message hub in \ours serves as a broadcast mechanism that simplifies group interactions. Developers can initiate a message hub by defining participating agents and can include initial broadcast messages. 
When new messages are generated by the agents within the message hub, they are automatically disseminated to other participants, as demonstrated in \lstlistingname~\ref{lst:msghub}. 
This abstraction is particularly useful for multi-agent scenarios involving LLMs, where dynamic and contextually rich conversations are commonly observed~\citep {du2023improving}.

\begin{lstlisting}[language=Python, float=h, caption={Using message hub with \ours.}, label={lst:msghub}]
# set up agents: agent1 to agent4
# ...

greeting = Msg("host", "Welcome to the message hub!")

with msghub(participant=[agent1, agent2, agent3], 
            announcement=greeting) as hub:
    # Message will be broadcast to agent2 and agent3 automatically
    agent1()
    
    # Delete agent2 from the message hub
    hub.delete(agent2)
    
    # Add agent4 into the message hub
    hub.add(agent4)
    
    # Broadcast message
    hub.broadcast(Msg("host", "Welcome agent4 to join the hub!"))
\end{lstlisting}  

\subsection{Resource-Rich Environment for Agent Development}
To further enhance usability, \ours is equipped with a rich set of built-in resources, including services, dedicated agents, and pre-configured examples. These resources are designed to reduce the initial setup effort and enable rapid prototyping and deployment of multi-agent LLM systems.

\paragraph{Comprehensive Service Integration}
\ours integrates various service functions, such as web search, database querying, and code execution, to support the tool usage capabilities of agents. 
These service functions are essential for building helpful agents with LLMs, as agents often need to draw information from external sources or execute tasks that go beyond the equipped LLMs' internal knowledge. \lstlistingname~\ref{lst:web_search} showcases the seamless conversion of a service into an OpenAI-Compatible JSON format, simplifying the integration process for developers.

\begin{lstlisting}[language=Python, float=t, caption={ Converting web search service into the function and JSON format dictionary that agent can use.}, label={lst:web_search}]
from agentscope.service import ServiceFactory, web_search

bing_search, func_json = ServiceFactory.get(web_search, engine="bing", api_key="xxx", num_results=10)

print(func_json)
# {
#     "name": "web_search",
#     "description": "Searching the given question with bing.",
#     "parameters": {
#         "type": "object",
#         "properties": {
#             "type": "object",
#             "properties": {
#                 "question": {
#                     "type": "string",
#                     "description": "The string question to search in Bing."
#                 }
#             }
#         }
#     }
# }

searching_result = bing_search("What's the date today?")
\end{lstlisting}

\paragraph{Pre-built Agent Templates}
As cataloged in Table \ref{tab:agents}, \ours offers pre-built agents and ready-to-use components for tasks like dialogue management, user proxying, multi-modal data handling, and distributed deployment. These templates serve as starting points for developers to customize and extend, significantly accelerating the development of multi-agent LLM applications.

\begin{table}[ht]
    \centering
    \begin{tabular}{l|l}
        \toprule
        Agent Name                 &    Function            \\ \hline
        UserAgent                  &    The proxy of the user.  \\
        DialogAgent                &    A general dialog agent, whose role can be set by system prompt.\\
        DictDialogAgent            &    A dictionary version dialog agent, who responds in Python dictionary format.  \\
        ReActAgent   & An agent that can reason and use tools \\
        ProgrammerAgent     &    An agent that can write and execute Python code.   \\
        TextToImageAgent       &    An agent that generates images according to the requirements.  \\
        % AudioDialogAgent         &    An agent that can interact in speech.   \\
        RpcUserAgent               &    A distributed version user proxy.   \\
        RpcDialogAgent             &    A distributed version DialogAgent.   \\
        \bottomrule
    \end{tabular}
    \caption{Some examples of built-in agents and their functions in \ours.}
    \label{tab:agents}
\end{table}

\subsection{Multi-Agent Oriented Demonstration Interfaces}
\label{sec:interactive}

Furthermore, \ours introduces interaction interfaces tailored for multi-agent systems, as illustrated in Figures \ref{fig:modal_terminal} and \ref{fig:modal_webui}. These interfaces provide a rich multi-modal experience, crucial for systems incorporating LLMs that handle diverse data types.


\begin{figure}[ht]
    \centering
    \includegraphics[width=0.75\linewidth]{figures/terminal.png}
    \caption{The dialogue history of a werewolf game in \ours.}
    \label{fig:modal_terminal}
\end{figure}

\begin{figure}[ht]
    \centering
    \fbox{\includegraphics[width=0.86\linewidth]{figures/web_ui_run.png}}
    \caption{Multi-modal interactions between agents in web UI.}
    \label{fig:modal_webui}
\end{figure}

\paragraph{Agent Differentiation in User Interfaces}
To facilitate user interaction with multiple agents, \ours assigns unique colors and icons to each agent, enhancing clarity and visual distinction in both terminal and web UI (\figref{fig:modal_webui}). The ``first-person perspective'' feature allows users to experience interactions from the viewpoint of a specified agent, aligning with their role in the application, such as in a game scenario. This feature not only enriches the multi-agent experience but also mirrors the nuanced interactions that occur in human-agent and agent-agent dialogues within LLM systems.

\paragraph{Monitoring and Cost Management}
A vital aspect of deploying LLMs at scale is resource management. \ours includes a monitoring module that tracks model and API usage, as well as calculating financial costs. Developers can customize metrics and set budget limits, receiving automatic alerts when thresholds are approached or exceeded. This proactive cost management is particularly important for LLMs that may incur high computational expenses.

\paragraph{AgentScope Gradio Interface}
Once you have a multi-agent application, executing it in the terminal may be a concise choice but lacks attraction.
In \ours, we provide a powerful Gradio-based interface that is compatible with all \ours applications as long as there is a \texttt{main} function as the application's entry point. 
For example, if the main function of the application is in \texttt{application.py} file, then running ``$\texttt{as\_studio application.py}$'' can build a Gradio application with a graphical user interface and support multi-modal content upload and presentation.


\subsection{Towards Graphical Application Development}

The design mentioned above provides massive convenience for those familiar with Python programming to quickly develop their multi-agent applications. However, \ours takes a step further. \ours provides a drag-and-drop online workstation on which developers only need to drag the module blocks to compose an application; then, the workstation can generate a configuration file of the application in JSON or even a piece of Python code. With this feature, those with limited experience with Python programming can build their multi-agent application without writing any Python code, while those familiar with Python can instantly obtain a piece of draft code ready for further customization.
A screenshot of the online workstation is shown in \figref{fig:workstation}, and the idea supporting this implementation is illustrated as follows.

\begin{figure}[ht]
    \centering
    \fbox{\includegraphics[width=0.86\linewidth]{figures/workstation.pdf}}
    \caption{Drag-and-drop programming \href{https://agentscope.aliyun.com/}{workstation}.}
    \label{fig:workstation}
\end{figure}

\paragraph{Expressing Multi-agent application with nodes in directed acyclic graph (DAG).}
Based on the highly modular design of our basic infrastructure, all the key components can be represented as a node, and an application can be built by constructing a directed cycle graph (DAG). The execution of the application is equivalent to triggering and running the nodes in the graph following the traversing order of DAG. Following the traditional terms, we name such DAG execution as a \emph{workflow} and name the nodes in the workflow as \emph{workflow nodes}.  
According to their functionality, the workflow nodes are categorized into six different types: model nodes, agent nodes, pipeline nodes, service nodes and copy nodes.  
\begin{itemize}
    \item 
    {\bf Model nodes}: Model nodes are designed to be relatively independent of the DAG. They correspond to the model configurations in \ours and work as entries to let users configure their models (LLMs, embedding models, or multi-modal models) and maintain such information for all the nodes in the following workflow that need to use the model.
    
    \item 
    {\bf Service (tool) nodes}: These nodes correspond to the services available in \ours. Some of them require additional information to set up, such as Google search and Bing search, which require API keys; others can be used directly. 
    
    \item 
    {\bf Agent nodes}: As the name suggests, agent nodes represent the agents in \ours, which means users need to decide the models, agent name and system prompts for the agent. 

    \item 
    {\bf Pipeline nodes}: The pipeline node includes the operators of \ours, including the message hub and the pipelines (sequential, for-loop, while-loop, etc.). With such nodes, DAG representations can be as concise as Python programming.
    
    \item 
    {\bf Message node}: The message node is designed for cases where some initial messages are needed, such as the announcement (initial message) for the message hub.
    
    \item 
    {\bf Copy node}: The copy node is a special kinds of node that replicate the results of a parent node when its output is needed for multiple subsequent operations.
\end{itemize}


\paragraph{Execute DAG with JSON or compile to Python.}
With the nodes above, developers can build applications by composing DAGs. However, the DAG is highly UI-dependent. Although a DAG can be represented in some formats (e.g., JSON format recording each node's information and execution dependency), we still need to ensure it is as reusable as other applications. To overcome this, \ours is equipped with a data structure called ASDiGraph, which provides two solutions based on it. 
\begin{itemize}
    \item 
    {\bf Direct-run}: Given a JSON file recording the DAG information, ASDiGraph can parse DAG information and sort the nodes in topological order. With these sorted nodes, the run function of ASDiGraph can execute them in order and feed the predecessor's output to their successors as an application is executed step by step.

    \item 
    {\bf To-Python compiler}: The second solution is to translate the JSON file to a Python script. With the highly modularized components of \ours, the key idea is to rely on internal mappings of the functionality, required inputs, and expected outputs to small pieces of Python code. Specifically, each node contains Python code for importing dependent modules, initiating models or agents, and executing the application logic. ASDiGraph first groups the pieces of importing code and initiating code, and then it composes the pieces of execution code following the topological order.
    Therefore, users will obtain a complete Python script after the ASDiGraph finishes compilation.
\end{itemize}





\subsection{Automatic Prompt Tuning}
\label{sebsec:prompt_tuning}
For a multi-agent system that utilizes LLMs for generation, writing an appropriate prompt requires significant human effort and expertise~\citep{pryzant2023automatic}, which motivates us to provide automatic prompt generation and tuning in \ours for its high usability. Specifically, \ours allows users to generate prompts based on a simple description of the agent in natural language, update prompts according to contexts, and enable in-context learning.

\paragraph{System Prompt Tuning} When an agent is created, a system prompt should be associated with the agent to define its roles and responsibilities for following human instructions. For example, a Programmer Agent might be prompted as ``You are proficient in writing and executing Python code''. Meanwhile, a detailed and informative system prompt can improve agent performance and ensure that the agent performs as expected, such as ``You are proficient in writing and executing Python code. You prefer to write the code in a modular fashion and provide unit tests for each module''. With \ours, users only need to provide a simple description of the agent when creating the agents, and \ours can automatically generate such helpful system prompts using built-in tools based on LLMs, as shown in Example~\ref{lst:prompt_generation}.
\begin{lstlisting}[language=Python, float=h, caption={Initialize a programmer agent with automatic system prompt generation.}, label={lst:prompt_generation}]
# set up agents with automatic prompt generation
# ...
from agentscope.agents import ProgrammerAgent

# Load model configs
agentscope.init(model_configs="model_configs.json")

# Create a programmer agent 
programmer_agent = ProgrammerAgent(name="assistant", auto_sys_prompt=True,
                        model_config_name="my_config",
                        sys_prompt="an assistant that can write Python code")
\end{lstlisting}

Besides, \ours provides interfaces for system prompt updates, which include manually setting by users or automatically adjusting based on the context. As a promising future direction, meta-prompting techniques~\citep{pryzant2023automatic, suzgun2024meta} can also be integrated into \ours, which might involve integrating an evaluator to provide guidance for automatic prompt optimization.

\paragraph{In-Context Learning} Providing multiple demonstrations to the LLMs can greatly enhance their ability to follow instructions, particularly when we want them to complete specific downstream tasks~\citep{dai2023gpt, wei2022emergent}. \ours provides a simple switch to turn on/off the in-context learning behavior for agents that utilize LLMs. When users choose to apply in-context learning, they only need to provide demonstration candidates and configure how to match the most suitable ones, as illustrated in Example~\ref{lst:icl}. \ours offers several widely-used and useful matching approaches, such as random selection, similar questions, and similar answers, and allows for user customization.

\begin{lstlisting}[language=Python, float=h, caption={Enable in-context learning when creating an agent.}, label={lst:icl}]
# set up agents with in-context learning
# ...
from agentscope.agents import ReActAgent
from agentscope.utils.common import load_demo_data

# Load model configs
agentscope.init(model_configs="model_configs.json")

# Load demonstrations
react_pairs = load_demo_data("my_demos.txt")

# Create a reAct agent 
react_agent = ReActAgent(name="react_agent", enable_icl=True, 
                         demos=react_pairs, matching_approach="random")
\end{lstlisting}